\section{Expanded Native30 evaluation including YuE2}
\label{sec:yue2-expanded}
The completed expanded development cohort contains 4,228 recordings: the
original 3,830 plus 398 eligible YuE2 development recordings. Another 100
YuE2 recordings remain a separate locked population, evaluated in the
prediction-only sensitivity analysis rather than inserted into this training
matrix. Two short originals are excluded from Native30. Shared Muse source-song
groups link generated recordings across models and are retained in the
group/component split restrictions.

All 255 subsets of S, D, R, P, F, H, SC and BC are evaluated under five
training-group caps and the three protocols. The completed evaluation comprises
112,455 model instances (80,325 primary; 32,130 diagnostic), with 3,570
structurally ineligible scheduled cells explicitly omitted. The reporting
projection contains 3,825 primary protocol/subset/cap cells. Training-only
transforms, fixed ridge parameter 10, threshold 0.5, null-value handling and
the group/source-balanced estimand follow the preceding Native30 protocol.

\begin{table}[htbp]\centering\small
\begin{tabular}{lrrr}\toprule
Feature set & Generator holdout & Human-source holdout & Grouped descriptive\\\midrule
S & 65.83 & 64.34 & 71.98\\
D & 59.23 & 63.42 & 63.58\\
R & 62.98 & 62.53 & 63.20\\
P & 53.32 & 54.42 & 54.78\\
F & 58.94 & 62.11 & 68.67\\
H & 54.69 & 59.51 & 62.40\\
SC & 56.07 & 62.63 & 64.76\\
BC & 46.11 & 45.59 & 48.02\\
S+D+R+P & 69.72 & 68.63 & 76.37\\
S+D+R+P+BC & 69.10 & 68.26 & 76.19\\
S+D+R+P+F+H+SC & 66.18 & 70.21 & 77.87\\
S+D+R+P+F+H+SC+BC & 66.29 & 70.25 & 78.37\\
\bottomrule\end{tabular}
\caption{Expanded Native30 group/source-balanced accuracy (\%) at the all-cap configuration. Protocols have distinct estimands; no cross-column causal comparison is implied.}
\label{tab:yue2-expanded-allcap}\end{table}
\begin{table}[htbp]\centering\small
\begin{tabular}{lrrrrr}\toprule
Protocol (all eight families) & 25 & 50 & 100 & 200 & All\\\midrule
Generator holdout & 65.46 & 64.11 & 64.48 & 65.94 & 66.29\\
Human-source holdout & 67.62 & 69.05 & 70.28 & 70.22 & 70.25\\
Grouped descriptive & 76.66 & 78.22 & 77.90 & 78.30 & 78.37\\
\bottomrule\end{tabular}
\caption{All-eight-family training-group-cap ablation. Caps are per-source group budgets, not total training-song counts. The complete 3,825-cell table retains all 255 subsets at all five caps.}
\label{tab:yue2-expanded-caps}\end{table}


At the all-cap configuration, S/D/R/P attains 69.7234\% balanced accuracy in
generator holdout, compared with 66.1827\% for seven families and 66.2877\%
for eight. In human-source holdout, the corresponding values are 68.6270\%,
70.2078\% and 70.2487\%; ordinary grouped descriptive values are 76.3652\%,
77.8734\% and 78.3713\%. Thus broader feature combinations have different
effects under different generalization targets. The results do not justify a
single universally superior configuration.

BC adds only 0.1050 percentage points to the seven-family generator-holdout
result and 0.0408 points to human-source holdout, versus 0.4980 points in the
ordinary grouped protocol. These small descriptive differences are not
significance evidence or proof that BC is a transferable authorship signature.
BC alone remains below 50\% in all three protocols at the fixed operating point.

\paragraph{Comparison with the pre-YuE2 cohort.}
The generator-holdout S/D/R/P result changes from 65.7872\% to 69.7234\%.
However, the expanded generator set changes both training composition and the
equally weighted held-generator test population. This is not a fixed-test-set
learning-curve estimate and cannot identify the causal benefit of adding
398 recordings. Similarly, altered component structure and missingness affect
eligible folds. The within-cohort cap tables and complete subset tables are
retained; choosing the best observed cap would require a further independent
selection/evaluation design.

\paragraph{Artifacts and evidence limits.}
The full report is retained under
\path{current_results/expanded_native30_report_v1/}. Its COMMIT SHA256 is:
\begin{quote}\scriptsize\ttfamily
b7934f2c649e5d0fbada832c335d36839a1aacb3bd7d6a674db83af5372bc911
\end{quote}
All seven committed report products were verified after local transfer.
The original model-evaluation COMMIT is:
\begin{quote}\scriptsize\ttfamily
11f074f867d0e3721c6a586e8c8f86116ff1b0198147ba14f6379b5d6f55570d
\end{quote}
\path{export_yue2_thesis_tables_v1.py} generates the complete 3,825-row
CSV and the displayed tables from that committed report, without refitting or
selection. The CSV resides in
\path{current_results/yue2_expanded_thesis_tables_v1/}. Reporting is not a
separately implemented audit of all aggregate numerical calculations, and no
confidence intervals or pooled cross-model AUC are claimed. Native60 transfer
is a separate duration-selected experiment and remains pending at this snapshot.
